% GENERATED FILE. Edit canonical lessons, then run npm run generate:books.
% canonical-source-hash: fnv1a64:32d6103355483d1a
% canonical-lessons: ZH-C501-yundongyuan, ZH-C501-guanjun, ZH-C501-jinpai, ZH-C501-huode, ZH-C501-qiuxing

\chapter{Winners}
\label{ch:zh-winners}
\hlchaptermodality{501}

\begin{chapteropening}
\textbf{By the end of this chapter:} \emph{I can talk about athletes, champions, gold medals, winning prizes and star players.}

Complete the last lesson of chapter 501: I can talk about athletes, champions, gold medals, winning prizes and star players.
\end{chapteropening}

\section[yùndòngyuán]{yùndòngyuán — athlete}

\label{lesson:ZH-C501-yundongyuan}



\begin{quote}
Before the new one: say the Mandarin for the score of a game, then the Mandarin for to take the field.
\end{quote}

\subsection*{You'll want to know: yùndòngyuán}
\emph{yùndòngyuán} — \textquotedblleft{}athlete\textquotedblright{}.

\begin{quote}
\emph{wǒ xiǎoshíhou xiǎng dāng yùndòngyuán} — When I was a kid I wanted to be an athlete.
\end{quote}

Its characters come later; for now it is a word for the ear.

\begin{cousinweb}
\emph{yùndòngyuán} is built from \emph{yùndòng} 'sport' and \emph{yuán} 'a member', all inherited from Middle Chinese, \emph{hjunH}, \emph{duwngX} and \emph{hjwen} (Guangyun, in Baxter's transcription), which glosses the first 'to move' and the third 'a count of things'.
\end{cousinweb}

\begin{grammarlens}[title={What you've built}]
One more word for saying how it was.
\end{grammarlens}

\subsection*{Your turn}
\begin{itemize}
\raggedright
  \item \emph{Say it:} \emph{yùndòngyuán}
  \item \emph{Say it:} \emph{yùndòngyuán}, once more
  \item \emph{Say it:} \emph{wǒ xiǎoshíhou xiǎng dāng yùndòngyuán} — When I was a kid I wanted to be an athlete.
  \item \emph{Recall:} say the Mandarin for the score of a game, then the Mandarin for to take the field, then say \emph{yùndòngyuán} again
\end{itemize}

\begin{tcolorbox}[breakable,colback=teal!4,colframe=teal!35!black,title={Before you move on}]
What does yùndòngyuán mean? (\textquotedblleft{}Athlete\textquotedblright{}.) Say it once more.
\end{tcolorbox}
\section[guànjūn]{guànjūn — champion}

\label{lesson:ZH-C501-guanjun}



\begin{quote}
Before the new one: say the Mandarin for to take the field, then the Mandarin for an athlete.
\end{quote}

\subsection*{You'll want to know: guànjūn}
\emph{guànjūn} — \textquotedblleft{}champion\textquotedblright{}.

\begin{quote}
\emph{tā dāng guo liǎng cì guànjūn} — He's been champion twice.
\end{quote}

Its characters come later; for now it is a word for the ear.

\begin{cousinweb}
\emph{guànjūn} is built from \emph{guàn} 'to put on a cap' and \emph{jūn} 'an army', both inherited from Middle Chinese, \emph{kwanH} and \emph{kjun} (Guangyun, in Baxter's transcription), which glosses the first 'to tie on the cap' and the second 'a body of troops'. How the pair came to mean 'champion' is not settled.
\end{cousinweb}

\begin{grammarlens}[title={What you've built}]
One more word for saying how it was.
\end{grammarlens}

\subsection*{Your turn}
\begin{itemize}
\raggedright
  \item \emph{Say it:} \emph{guànjūn}
  \item \emph{Say it:} \emph{guànjūn}, once more
  \item \emph{Say it:} \emph{tā dāng guo liǎng cì guànjūn} — He's been champion twice.
  \item \emph{Recall:} say the Mandarin for to take the field, then the Mandarin for an athlete, then say \emph{guànjūn} again
\end{itemize}

\begin{tcolorbox}[breakable,colback=teal!4,colframe=teal!35!black,title={Before you move on}]
What does guànjūn mean? (\textquotedblleft{}Champion\textquotedblright{}.) Say it once more.
\end{tcolorbox}
\section[jīnpái]{jīnpái — gold medal}

\label{lesson:ZH-C501-jinpai}



\begin{quote}
Before the new one: say the Mandarin for an athlete, then the Mandarin for the champion.
\end{quote}

\subsection*{You'll want to know: jīnpái}
\emph{jīnpái} — \textquotedblleft{}gold medal\textquotedblright{}.

\begin{quote}
\emph{tā zuì xiǎng yào yí kuài jīnpái} — What he wants most is a gold medal.
\end{quote}

Its characters come later; for now it is a word for the ear.

\begin{cousinweb}
\emph{jīnpái} is built from \emph{jīn} 'gold' and \emph{pái} 'a tablet', both inherited from Middle Chinese, \emph{kim} and \emph{bea} (Guangyun, in Baxter's transcription), which glosses the first 'gold' and the second 'a signboard'.
\end{cousinweb}

\begin{grammarlens}[title={What you've built}]
One more word for saying how it was.
\end{grammarlens}

\subsection*{Your turn}
\begin{itemize}
\raggedright
  \item \emph{Say it:} \emph{jīnpái}
  \item \emph{Say it:} \emph{jīnpái}, once more
  \item \emph{Say it:} \emph{tā zuì xiǎng yào yí kuài jīnpái} — What he wants most is a gold medal.
  \item \emph{Recall:} say the Mandarin for an athlete, then the Mandarin for the champion, then say \emph{jīnpái} again
\end{itemize}

\begin{tcolorbox}[breakable,colback=teal!4,colframe=teal!35!black,title={Before you move on}]
What does jīnpái mean? (\textquotedblleft{}Gold medal\textquotedblright{}.) Say it once more.
\end{tcolorbox}
\section[huòdé]{huòdé — to win (a prize); to be awarded}

\label{lesson:ZH-C501-huode}



\begin{quote}
Before the new one: say the Mandarin for the champion, then the Mandarin for a gold medal.
\end{quote}

\subsection*{You'll want to know: huòdé}
\emph{huòdé} — \textquotedblleft{}to win (a prize); to be awarded\textquotedblright{}.

\begin{quote}
\emph{tā huòdé le liǎng kuài jīnpái} — She won two gold medals.
\end{quote}

Its characters come later; for now it is a word for the ear.

\begin{cousinweb}
\emph{huòdé} is built from \emph{huò} 'to get' and \zh{得} 'to gain', both inherited from Middle Chinese, \emph{hweak} and \emph{tok} (Guangyun, in Baxter's transcription), which glosses the first 'to get' and the second 'gain and loss'.
\end{cousinweb}

\begin{grammarlens}[title={What you've built}]
One more word for saying how it was.
\end{grammarlens}

\subsection*{Your turn}
\begin{itemize}
\raggedright
  \item \emph{Say it:} \emph{huòdé}
  \item \emph{Say it:} \emph{huòdé}, once more
  \item \emph{Say it:} \emph{tā huòdé le liǎng kuài jīnpái} — She won two gold medals.
  \item \emph{Recall:} say the Mandarin for the champion, then the Mandarin for a gold medal, then say \emph{huòdé} again
\end{itemize}

\begin{tcolorbox}[breakable,colback=teal!4,colframe=teal!35!black,title={Before you move on}]
What does huòdé mean? (\textquotedblleft{}To win (a prize); to be awarded\textquotedblright{}.) Say it once more.
\end{tcolorbox}
\section[\texorpdfstring{\zh{球星} (qiúxīng)}{球星 (qiúxīng)}]{\zh{球星} (qiúxīng) — star player; sports star}

\label{lesson:ZH-C501-qiuxing}



\begin{quote}
Before the new one: say the Mandarin for a gold medal, then the Mandarin for to win (a prize).
\end{quote}

\subsection*{You'll want to know: \zh{球星}}
\textbf{\zh{球星}} — \emph{qiúxīng} — \textquotedblleft{}star player; sports star\textquotedblright{}.

\begin{quote}
\textbf{\zh{那个球星今天没有来比赛}} \emph{nàge qiúxīng jīntiān méiyǒu lái bǐsài} — That star player didn't come to play today.
\end{quote}

\begin{cousinweb}
\textbf{\zh{球星}} (\emph{qiúxīng}) is built from \zh{球} 'a ball' and \zh{星} 'a star', both inherited from Middle Chinese, \emph{gjuw} and \emph{seng} (Guangyun, in Baxter's transcription), which glosses the second 'the stars'.
\end{cousinweb}

\begin{grammarlens}[title={What you've built}]
One more word for saying how it was.
\end{grammarlens}

\subsection*{Your turn}
\begin{itemize}
\raggedright
  \item \emph{Say it:} \emph{qiúxīng}
  \item \emph{Say it:} \emph{qiúxīng}, once more
  \item \emph{Say it:} \emph{nàge qiúxīng jīntiān méiyǒu lái bǐsài} — That star player didn't come to play today.
  \item \emph{Recall:} say the Mandarin for a gold medal, then the Mandarin for to win (a prize), then say \emph{qiúxīng} again
\end{itemize}

\begin{tcolorbox}[breakable,colback=teal!4,colframe=teal!35!black,title={Before you move on}]
What does \zh{球星} mean? (\textquotedblleft{}Star player; sports star\textquotedblright{}.) Say it once more.
\end{tcolorbox}
